% Source: https://artofproblemsolving.com/community/c3h2211085_descartes_circle_forumla
% Stable identity: aops-topic-2211085-q01
\textbf{(a)} The centres of the three unit circles are at distance $2$ from each other, so they form an equilateral triangle of side $2$. Its centre $O$ is at distance $\dfrac{2}{\sqrt3}$ (the circumradius) from each vertex. By symmetry both required circles are centred at $O$.

\begin{asy}
usepackage("amsmath");
size(6cm);
defaultpen(fontsize(10pt));
real s=2/sqrt(3);
pair O=(0,0);
pair P1=s*dir(90), P2=s*dir(210), P3=s*dir(330);
draw(circle(O,s+1),blue+linewidth(0.9));
draw(circle(P1,1)); draw(circle(P2,1)); draw(circle(P3,1));
draw(circle(O,s-1),red+linewidth(0.9));
draw(P1--P2--P3--cycle,dashed+gray);
dot(P1);dot(P2);dot(P3);dot(O);
draw(P2--(P2+dir(210)),Arrow(4)); label("$1$",P2+0.5*dir(210),NW);
draw(O--(O+(s+1)*dir(20)),blue,Arrow(4)); label("$R$",(s+1)*dir(20)*0.75,N,blue);
\end{asy}

The small circle touches each unit circle from outside, so its radius is the distance from $O$ to a centre minus $1$:
\[r=\dfrac{2}{\sqrt3}-1=\boxed{\dfrac{2\sqrt3-3}{3}}\approx0.155.\]
The large circle touches each unit circle from outside it (internally), so its radius is that distance plus $1$:
\[R=\dfrac{2}{\sqrt3}+1=\boxed{\dfrac{2\sqrt3+3}{3}}\approx2.155.\]

\textbf{(b)} With $k_1=k_2=k_3=1$ the equation becomes $(3+k)^2=2(3+k^2)$, i.e. $k^2-6k-3=0$, so
\[k=3\pm2\sqrt3.\]
\begin{itemize}
\item $k=3+2\sqrt3>0$: then $\dfrac{1}{k}=\dfrac{1}{2\sqrt3+3}=\dfrac{2\sqrt3-3}{12-9}=\dfrac{2\sqrt3-3}{3}=r$. This is the small circle.
\item $k=3-2\sqrt3<0$: then $-\dfrac{1}{k}=\dfrac{1}{2\sqrt3-3}=\dfrac{2\sqrt3+3}{3}=R$. This is the enclosing circle, with curvature $-\dfrac{1}{R}$.
\end{itemize}
If instead we used $k_4=+\dfrac{1}{R}=2\sqrt3-3$, the left side would be $(3+2\sqrt3-3)^2=12$, while the right side would be $2\left(3+(2\sqrt3-3)^2\right)=2(24-12\sqrt3)=48-24\sqrt3\approx6.43$, so the formula fails.

The sign records the direction of tangency: the other circles touch the enclosing circle from the inside, where it bends around them instead of away from them, so its curvature counts as negative.
